%%%%%%%%%%  scan idx 207  |  folio 194  %%%%%%%%%%
ce qui achève la
démonstration de a).

\emph{Preuve de} b). Le cohéreur, étant de dimension cohomologique finie d'après a), admet ([H] I 5.3) un dérivé droit
\[
 \mathbb RQ:\mathrm D(S)\longrightarrow\mathrm D(\Qcoh(S))
\]
(dont la restriction à $\mathrm D^+(S)$ coïncide avec le foncteur $\mathbb RQ$ défini plus haut). Pour $E\in\Ob\mathrm D(\Qcoh(S))$, il existe un isomorphisme fonctoriel canonique
\begin{equation}\tag{1}
 E\xrightarrow{\sim}\mathbb RQ(E),
\end{equation}
puisque, d'après (3.5.1), les $E^i$ sont acycliques pour $Q$ et tels que $E^i\xrightarrow{\sim}Q(E^i)$. D'autre part, pour $E\in\Ob\mathrm D(S)$, on a un homomorphisme fonctoriel canonique
\begin{equation}\tag{2}
 \mathbb RQ(E)\longrightarrow E
\end{equation}
(provenant, dans le cas où les $E^i$ sont $Q$-acycliques, de l'homomorphisme d'adjonction ordinaire entre $\varphi$ et $Q$). On vérifie trivialement que les homomorphismes (1) et (2) font de $\mathbb RQ:\mathrm D(S)\longrightarrow\mathrm D(\Qcoh(S))$ un adjoint à droite de $\varphi:\mathrm D(\Qcoh(S))\longrightarrow\mathrm D(S)$. Le fait que (1) soit un isomorphisme implique que $\varphi$ est pleinement fidèle. Enfin, utilisant la suite spectrale
\[
 R^pQ(H^q(E))\Longrightarrow R^{*}Q(E),
\]
(convergente puisque $Q$ est de dimension cohomologique finie), on voit que (2) est un isomorphisme pour $E\in\Ob\mathrm D(S)_{\qcoh}$, donc que $\varphi$ a pour image essentielle $\mathrm D(S)_{\qcoh}$. Cela prouve b) et achève la démonstration de (3.7).

